\documentclass[11pt]{article}
\usepackage[margin=2.5cm]{geometry}
\usepackage{graphicx,booktabs,hyperref,amsmath,amssymb}
\usepackage[T1]{fontenc}

\title{Cross-Model CKA Alignment Between EVO2 and RFdiffusion3\\
Reflects Low-Frequency Smoothing Bias\\
Rather Than Structural Convergence}

\author{Cheng Ding\\
\textit{ShanghaiTech University}\\
\texttt{dingcheng2024@shanghaitech.edu.cn}}

\date{}

\begin{document}
\maketitle

\begin{abstract}
We test whether DNA and protein language models (EVO2, FAESM/ESM2-650M)
share intermediate head-level residue representations with a
protein-structure diffusion model (RFdiffusion3) on the AsCas12f1
CRISPR-Cas locus. After correcting the protein length to 422 aa
(UniProt A0A2U3D0N8) and aligning coordinates via codon pooling,
we scan 160 EVO2 heads, 660 FAESM heads, and 1,216 RF3 heads
(194,560--802,560 pairs). The strongest raw head pair achieves
CKA~=~0.851 (2,000-permutation max-null $p \approx 0.0005$).

However, extensive controls reveal that high raw CKA is dominated by
multiscale residue-order geometry rather than domain-specific functional
alignment. Position-only baselines match or exceed real EVO2 at segment
scales (mean excess $-$0.43 to $-$0.48). Apparent domain enrichment
vanishes under position-aware nulls. Mutation-induced delta CKA is
indistinguishable from pair-shuffled controls ($p = 0.42$), and
retrieval of matched mutants is at chance. A protein-LM control
using FAESM reproduces the same pattern: high static CKA without
mutation-specific alignment.

The triangular EVO2--FAESM--RF3 comparison shows that cross-model CKA
primarily measures shared residue-order geometry rather than biological
semantics. High static CKA---whether from a DNA LM or a protein LM---does
not imply recoverable mutation-specific or residue-level semantic alignment.
\end{abstract}

\end{document}
